\section{The actual return and its multiplicity}\label{sec:return-framework}
A positive function is called a positive unit when it and its reciprocal
are bounded. Phase and section derivatives keep the field parameters
fixed unless a parameter derivative is specified.
Fix $J\Subset(0,2)$ and write $u=u_{\rm up}$ for the reserved
physical coordinate on the upper section. On the fast arc choose
\[
 \Sigma_u^b=\{(0,y_b+u):|u|<u_0\},\qquad
 y_b=1+\exp(\pi b/\sqrt{4-b^2}),
\]
and a small outgoing section $\Sigma_\xi$ at the lower intersection
with $z=\xi>0$. The source and target cores are fixed before the
parameters are shrunk. Let $T_-$ be the backward regular fast transition
and $T_+$ the forward transition through descent, the slow segment,
and contact. The displacement
\begin{equation}\label{eq:common-displacement}
 \Delta(u)=T_+(u)-T_-(u)
\end{equation}
is defined only where both passages exist in finite time.
The target uses the fast first integral $E_b$ of
Lemma~\ref{lem:fast-energy}. In these orientations
$T_-'=-j$, with $j$ and its logarithmic derivatives bounded on the
fixed regular core.

Lemma~\ref{lem:sections} proves that every whole-graphic cycle lies in
$y>0$ and meets each prescribed regular section exactly once.
One Hausdorff inclusion places the orbit near every graphic marker,
forcing a crossing inside the prescribed core; the other confines all
crossings to that core. Jordan separation and the fixed transverse sign
then give uniqueness. The regular holonomies are defined for all initial
values in a smaller core, independently of periodicity.


\subsection{Pullback to actual minimal-return multiplicity}
For two physical sections of the same finite orbit, regular
holonomy $h$ gives $P_2=h\circ P_1\circ h^{-1}$.
Near a return zero,
\[
 h(P_1(s))-h(s)
 =(P_1(s)-s)\int_0^1h'(s+t(P_1(s)-s))\,dt .
\]
The last factor is nonzero at the zero.
A signed energy target or a scaled section change has the same
property on its actual finite-flow domain.
Uniform invertibility at a singular parameter limit is unnecessary;
fixed-field nonzero derivatives suffice.
The regional estimates detect simple zeros or stationary zeros with
nonzero second derivative. Their finite-order section coordinates
therefore determine the order of the physical analytic return zero.

Equilibria and infinite saddle passages bound the actual domain.
At a finite event, two passage formulas describe the same flow at one
source value. Each simple displacement zero has actual multiplicity one,
and each stationary zero with nonzero second derivative has multiplicity two.

\subsection{Normalized coefficient directions}
For each $\mu\ne0$ define $\varepsilon>0$ by
\begin{equation}\label{eq:weighted-radius-main}
 \frac{\mu_1^2+\mu_2^2+\mu_3^2}{\varepsilon^4}
                  +\frac{\mu_4^2}{\varepsilon^6}=1 .
\end{equation}
The left side strictly decreases from infinity to zero, so the
radius is unique. Put $E_i=\mu_{i+1}/\varepsilon^2$ and
$D=\mu_4/\varepsilon^3$. This is a re-expression of the same field;
it does not alter phase coordinates or time. Its direction belongs
to $S^3$, and $\varepsilon(\mu)\to0$ as $\mu\to0$.

Write, exactly,
\begin{equation}\label{eq:upper-normalized}
 F=y(1+bz-y)+\varepsilon^2P(z),\qquad
 G=zy+\varepsilon^3D,\qquad
 P(z)=E_0+E_1z+E_2z^2 .
\end{equation}
All regional estimates below hold on open coefficient windows with
fixed finite phase order, fixed section cores, and one whole-graphic
phase neighborhood.

